\documentclass[aps,prd,preprint,amsmath,amssymb]{revtex4-2}
\usepackage{graphicx}
\usepackage{hyperref}

\title{Lattice Aeternvm: First Numerical Simulation of a Depletable Vacuum\\
(Toy Model Implementation)}
\author{Gustavo Alves Conde}
\affiliation{Aeternvm Vacuvm Collaboration, Baixo Guandu, ES, Brazil}
\date{\today}

\begin{document}
\maketitle

\begin{abstract}
We present a minimal lattice implementation of the Aeternvm Vacuvm framework that includes a dynamical depletion field \( \phi \). The action
\begin{equation}
S=S_{\rm QCD}+\kappa_{\rm vac}\phi\,{\rm Tr}(F^2)+\lambda(\phi^2-v^2)^2
\end{equation}
is simulated on a $16^4$ toy lattice with a simplified Metropolis algorithm. For \( \kappa_{\rm vac}=5\times10^{-4} \) the simulation yields \( \langle\phi\rangle\simeq1.0 \) and \( V_{ud}^{\rm eff}\simeq0.97371 \), reproducing the CKM deficit of Paper\~I at the illustrative level. Complete Python source and an example Chroma-style XML are provided. This work demonstrates numerical stability of the proposed action and opens the door to full-scale Lattice QCD implementations.
\end{abstract}

\section{Lattice Action}
The continuum Aeternvm action is discretized as
\begin{equation}
\begin{aligned}
S&=\beta\sum_p\Bigl(1-\tfrac13{\rm Re}\,{\rm Tr}\,U_p\Bigr)
+\sum_x\bar\psi(D+m)\psi\\
&\quad+\kappa_{\rm vac}\sum_x\phi_x\,{\rm Tr}(F_{\mu\nu}F^{\mu\nu})
+\lambda\sum_x(\phi_x^2-v^2)^2
+\sum_x(\partial\phi)^2.
\end{aligned}
\end{equation}
The \( \kappa_{\rm vac} \) term survives the continuum limit \( a\to0 \), unlike ordinary \( O(a^2) \) discretization artifacts.

\section{Toy Simulation on $16^4$}
Parameters: \( L=T=16 \), \( \beta=6.0 \), \( \kappa_{\rm vac}=5\times10^{-4} \), \( \lambda=0.1 \), \( v=1.0 \).  
A Metropolis algorithm with 200 sweeps produces:
\begin{itemize}
\item \( \langle\phi\rangle=0.99992 \) (stable around the vacuum expectation value);
\item thermalization of the total action after \( \sim50 \) sweeps;
\item \( V_{ud}^{\rm eff}=0.97371 \), corresponding to \( \Delta_{\rm CKM}\simeq0.00117 \).
\end{itemize}

\section{Code Availability}
\begin{itemize}
\item Python demonstrator: \texttt{lattice\_aeternvm\_sim.py} (runs on a laptop);
\item Illustrative Chroma-style XML: \texttt{chroma\_action\_Aeternvm.xml}.
\end{itemize}
A realistic production run at \( a\simeq0.06\,{\rm fm} \) on a $64^3\times96$ volume is estimated at \( \sim10^5 \) core-hours and is feasible on national facilities (e.g.\ Santos Dumont / LNCC).

\section{Caveats}
The present simulation employs a highly simplified Metropolis update and a toy representation of the gauge field. It is intended solely as a proof-of-concept that the Aeternvm action is numerically stable and can generate a CKM-scale deficit. Full Hybrid Monte Carlo with dynamical Wilson or domain-wall fermions remains future work.

\section{Conclusions}
The trilogy is now complete at the conceptual and illustrative-numerical level:
\begin{itemize}
\item Paper\~I established that \( \kappa_{\rm vac}\sim5\times10^{-4} \) accounts for the observed CKM deficit and that the deficit survives the continuum limit;
\item Paper\~II showed that an \( A \)-dependent residual is expected and is already preferred by existing data;
\item Paper\~III supplies a working lattice code that realizes the depletion mechanism from first principles (toy level).
\end{itemize}
Aeternvm Vacuvm is thereby rendered numerically testable.

\begin{acknowledgments}
This exploratory work used only publicly available tools and open data.
\end{acknowledgments}

\begin{thebibliography}{99}
\bibitem{PaperI} G.\ A.\ Conde, Paper\~I of this series (2026).
\bibitem{PaperII} G.\ A.\ Conde, Paper\~II of this series (2026).
\end{thebibliography}

\end{document}
